\subsection*{Procedimiento y operaciones del simulacro S07}
Se emplean el círculo de Mohr, la forma de Larson-Miller con logaritmo decimal, la ley de Arrhenius para el coeficiente de difusión y la relación entre tiempo de vuelo y profundidad $d=vt/2$.

\begin{enumerate}

\item Tensión normal en x.

\[\fitmath{\mathrm{sx}=120\approx 120\,\mathrm{MPa}}\]

\item Tensión normal en y.

\[\fitmath{\mathrm{sy}=-40\approx -40\,\mathrm{MPa}}\]

\item Tensión tangencial.

\[\fitmath{\mathrm{txy}=60\approx 60\,\mathrm{MPa}}\]

\item Centro del círculo de Mohr.

\[\fitmath{c=\frac{\mathrm{sx}+\mathrm{sy}}{2}}\]

\[\fitmath{c=\frac{(120)+(-40)}{2}\approx 40\,\mathrm{MPa}}\]

\item Radio del círculo de Mohr.

\[\fitmath{\mathrm{rad}=\sqrt{{\left(\frac{\mathrm{sx}-\mathrm{sy}}{2}\right)}^{2}+{\mathrm{txy}}^{2}}}\]

\[\fitmath{\mathrm{rad}=\sqrt{{\left(\frac{(120)-(-40)}{2}\right)}^{2}+{(60)}^{2}}\approx 100\,\mathrm{MPa}}\]

\item Tensión principal mayor.

\[\fitmath{s_{1}=c+\mathrm{rad}}\]

\[\fitmath{s_{1}=(40)+(100)\approx 140\,\mathrm{MPa}}\]

\item Tensión principal menor.

\[\fitmath{s_{2}=c-\mathrm{rad}}\]

\[\fitmath{s_{2}=(40)-(100)\approx -60\,\mathrm{MPa}}\]

\item Tensión tangencial máxima en el plano.

\[\fitmath{\mathrm{tmax}=\mathrm{rad}}\]

\[\fitmath{\mathrm{tmax}=(100)\approx 100\,\mathrm{MPa}}\]

\item Orientación del plano principal respecto al eje x.

\[\fitmath{\mathrm{th}=\frac{0{,}5\cdot \arctan\left(\frac{2\cdot \mathrm{txy}}{\mathrm{sx}-\mathrm{sy}}\right)\cdot 180}{\pi}}\]

\[\fitmath{\mathrm{th}=\frac{0{,}5\cdot \arctan\left(\frac{2\cdot (60)}{(120)-(-40)}\right)\cdot 180}{\pi}\approx 18{,}43495\,{}^{\circ}}\]

\item Tensión equivalente de von Mises en tensión plana.

\[\fitmath{\mathrm{svm}=\sqrt{{s_{1}}^{2}-s_{1}\cdot s_{2}+{s_{2}}^{2}}}\]

\[\fitmath{\mathrm{svm}=\sqrt{{(140)}^{2}-(140)\cdot (-60)+{(-60)}^{2}}\approx 177{,}7639\,\mathrm{MPa}}\]

\item Coeficiente de seguridad frente a un límite elástico de 350 MPa.

\[\fitmath{\mathrm{SF}=\frac{350}{\mathrm{svm}}}\]

\[\fitmath{\mathrm{SF}=\frac{350}{(177{,}7638883)}\approx 1{,}968904\,}\]

\item Temperatura de referencia.

\[\fitmath{T_{1}=1000\approx 1000\,\mathrm{K}}\]

\item Tiempo de rotura de referencia.

\[\fitmath{t_{1}=100\approx 100\,\mathrm{h}}\]

\item Constante de Larson-Miller.

\[\fitmath{C=20\approx 20\,}\]

\item Parámetro de Larson-Miller con tiempo en horas.

\[\fitmath{\mathrm{LMP}=T_{1}\cdot \left(C+\frac{\ln\left(t_{1}\right)}{\ln\left(10\right)}\right)}\]

\[\fitmath{\mathrm{LMP}=(1000)\cdot \left((20)+\frac{\ln\left((100)\right)}{\ln\left(10\right)}\right)\approx 22000\,\mathrm{K}}\]

\item Nueva temperatura.

\[\fitmath{T_{2}=1050\approx 1050\,\mathrm{K}}\]

\item Tiempo de rotura a la nueva temperatura con el mismo parámetro.

\[\fitmath{t_{2}=\exp\left(\left(\frac{\mathrm{LMP}}{T_{2}}-C\right)\cdot \ln\left(10\right)\right)}\]

\[\fitmath{t_{2}=\exp\left(\left(\frac{(22000)}{(1050)}-(20)\right)\cdot \ln\left(10\right)\right)\approx 8{,}961505\,\mathrm{h}}\]

\item Factor preexponencial de difusión del carbono en austenita.

\[\fitmath{D_{0}=2{,}3\times10^{-5}\approx 2{,}3\times10^{-5}\,\mathrm{m^{2}/s}}\]

\item Energía de activación.

\[\fitmath{Q=148000\approx 148000\,\mathrm{J/mol}}\]

\item Temperatura de cementación.

\[\fitmath{\mathrm{Tc}=1223\approx 1223\,\mathrm{K}}\]

\item Coeficiente de difusión de Arrhenius.

\[\fitmath{D=D_{0}\cdot \exp\left(\frac{-Q}{8{,}314462618\cdot \mathrm{Tc}}\right)}\]

\[\fitmath{D=(2{,}3\times10^{-5})\cdot \exp\left(\frac{-(148000)}{8{,}314462618\cdot (1223)}\right)\approx 1{,}098332\times10^{-11}\,\mathrm{m^{2}/s}}\]

\item Tiempo de tratamiento.

\[\fitmath{\mathrm{tc}=4\cdot 3600\approx 14400\,\mathrm{s}}\]

\item Profundidad de difusión característica.

\[\fitmath{x=\sqrt{D\cdot \mathrm{tc}}}\]

\[\fitmath{x=\sqrt{(1{,}098332263\times10^{-11})\cdot (14400)}\approx 0{,}0003976932\,\mathrm{m}}\]

\item Velocidad longitudinal del sonido en el aluminio, dato docente.

\[\fitmath{v=6300\approx 6300\,\mathrm{m/s}}\]

\item Tiempo de ida y vuelta en una pared de 25 mm.

\[\fitmath{\mathrm{tf}=\frac{2\cdot 0{,}025}{v}\cdot 1000000}\]

\[\fitmath{\mathrm{tf}=\frac{2\cdot 0{,}025}{(6300)}\cdot 1000000\approx 7{,}936508\,\mathrm{us}}\]

\item Profundidad de un eco recibido a 4{,}76 microsegundos.

\[\fitmath{\mathrm{de}=\frac{v\cdot 4{,}76\times10^{-6}}{2}\cdot 1000}\]

\[\fitmath{\mathrm{de}=\frac{(6300)\cdot 4{,}76\times10^{-6}}{2}\cdot 1000\approx 14{,}994\,\mathrm{mm}}\]

\item Diferencia de amplitud para una razón de 3.

\[\fitmath{\mathrm{dB}=\frac{20\cdot \ln\left(3\right)}{\ln\left(10\right)}\approx 9{,}542425\,\mathrm{dB}}\]

\end{enumerate}
